\section{Fase 3 - PatchAgent: riparazione strutturale e validazione}
\label{sec:fase3-validazione}

\subsection{Riparazione Strutturale via Tree-sitter: CST vs AST}
\label{subsec:3.4.5}
L'approccio alla modifica del codice sorgente rappresenta uno degli snodi 
architetturali più delicati nella progettazione diun sistema di Automated 
Program Repair \cite{huang2024evolving}. Tradurre un'intenzione correttiva 
in una mutazione fisica del file richiede un meccanismo che garantisca non 
solo l'accuratezza della modifica, ma anche la validità formale dell'artefatto 
risultante. Per soddisfare i severi requisiti di robustezza di\texttt{HybridRepair}, 
l'architettura adotta la manipolazione diretta del \textit{ConcreteSyntax Tree}
avvalendosi della libreria Tree-sitter, distaccandosi nettamente dai paradigmi testuali 
classici.
\subsubsection{I Limiti Fondamentali degli Approcci Testuali}
\label{subsubsec:3.4.5.1}
Prima di esplorare le specificità tecniche di Tree-sitter,è imperativo 
analizzare criticamente le ragioni per cuigli approcci tradizionali, 
manipolazione tramite espressioni regolari(regex), diff unificati e operatori di mutazione, 
risultano inadeguati per un sistema APR di alta affidabilità.

\paragraph{Il fallimento delle espressioni regolari e il problema del bilanciamento}\mbox{}\\
Le espressioni regolari operano analizzando flussi monodimensionali di caratteri. 
Questo approccio si scontra frontalmente con la natura gerarchica e 
ricorsiva del codice sorgente Java. Un costrutto come un metodo non è una semplice sequenza lineare,
 ma una struttura annidata che può includere blocchi condizionali \texttt{if/else}, cicli 
 \texttt{for/while}, blocchi \texttt{try/catch/finally}, oltre a costrutti complessi come 
 classi anonime o espressioni lambda. Ciascuno di questi elementi introduce nuovi livelli di 
 parentesi graffe che devono essere rigorosamente bilanciati. Dal punto di vista della teoria dei 
 linguaggi formali, il problema del bilanciamento delle parentesi appartiene alla classe dei linguaggi
  liberi dal contesto (\textit{context-free}), mentre le espressioni regolari possono descrivere
   unicamente linguaggi regolari \cite{aho2006compilers}. Una regex semplicistica 
   (come \verb|\w+\s*$[^)]*$\s*{[^}]*}|) è ingrado di catturare solo metodi piatti, 
   fallendo inesorabilmente di fronte al primo blocco condizionale annidato. Storicamente,
   la prima iterazione architetturale (v1) di \texttt{HybridRepair} tentava di arginare questa 
   limitazione combinando l'uso di regex con contatori manuali di parentesi graffe (\texttt{\_count\_braces()}). 
   Tuttavia, questo approccio euristico scatenava il cosiddetto ``problema P1'':la generazione di patch contenenti 
   parentesi sbilanciate o dichiarazioni di metodi frammentate o duplicate, costringendo il sistema a scartare un 
   elevato numero di tentativi validi.

\paragraph{Il fallimento dei diff unificati nei sistemi multi-tentativo}
Il formato \textit{unidiff} (standardizzato in strumenti come \texttt{git diff} e adottato 
da vecchi sistemi APR) codifica una modifica definendo un blocco testuale (l'\textit{hunk}) 
ancorato a specifiche righe di contesto preesistenti. Questo modello si fonda su un'assunzione
 di immutabilità del contesto: le righe circostanti la modifica devono coincidere in modo esatto
  con quelle del file bersaglio. In un agente iterativo guidato da un LLM, questa assunzione crolla miseramente. 
  L'agente esplora molteplici variazioni, e i diff generati dal modello spesso fanno riferimento al codice 
  sorgente fallato originale (fornito nel prompt iniziale). Tuttavia, al procedere dei round esplorativi,
  il file bersaglio può accumulare piccole variazioni o trovarsi in uno stato transitorio alterato. 
  Il disallineamento posizionale che ne consegue causa il rigetto del diff in fase di applicazione, 
  scaturendo in un falso negativo: la logica della patch potrebbe essere brillante e corretta, ma il 
  suo veicolo testuale (il diff) ne impedisce la materializzazione \cite{jimenez2024sweagent, xia2023chatrepair}.

  \paragraph{Il fallimento dei mutanti e l'overfitting}L'ultimo approccio tradizionale è quello basato sui 
  mutanti. Questi sistemi alterano l'AST manipolando schemi predefiniti 
  (operatori di mutazione), ad esempio forzando una condizione a \texttt{true}, sostituendo un operatore relazionale
   o inserendo controlli di nullità (\textit{null-checks}). Questa strategia soffre di un doppio limite. 
   In primo luogo, lo spazio esplorativo è confinato al ristretto set degli operatori implementati: se la 
   risoluzione del bug richiede la stesura di un'istruzione logica non codificata da un mutante predefinito, 
   la riparazione è matematicamente impossibile. In secondo luogo, e in maniera ancora più grave, i mutanti 
   inducono un massiccio \textit{overfitting} (sovradattamento). Le mutazioni casuali tendono a generare 
   patch che bypassano il caso di test fallito in modo parziale o accidentale, alterando o degradando 
   silenziosamente comportamenti del programma per input non coperti dai test. Ricerche empiriche dimostrano 
   che il tasso di \textit{overfitting} per alcuni di questi sistemi può superare il 90\% delle patch giudicate plausibili \cite{qi2015analysis, huang2024evolving}.

\subsubsection{Introduzione a Tree-sitter: il Parser Incrementale ad Alta Fedeltà}
\label{subsubsec:3.4.5.2}
Per superare le fragilità  degli approcci testuali,
\texttt{HybridRepair} adotta Tree-sitter, una libreria di parsing incrementale
e open-source sviluppata originariamente da GitHub. A differenza
dei parser canonici integrati nei compilatori Java (come
l'infrastruttura ECJ o \texttt{javac}), Tree-sitter è spinto da
una grammatica GLR (\textit{Generalized Left-to-Right Rightmost}).

Il vantaggio predominante di un parser GLR risiede
nella sua resilienza. Nei tradizionali parser di classe
LL(k) o LR(k), l'incontro di una deviazione sintattica
innesca un errore fatale, causando l'interruzione immediata dell'analisi
e la distruzione dell'albero parziale. Tree-sitter, invece, gestisce
le anomalie sintattiche incapsulandole all'interno di speciali nodi
\texttt{ERROR} nel suo albero, portando a termine il
parsing dell'intero documento. Questa capacità è di importanza
critica nel dominio dell'APR, ove il file sorgente
sottoposto ad analisi è, per definizione, difettoso e
potrebbe già presentare lievi imperfezioni sintattiche.

Dal punto di vista dell'ingegneria del software, l'integrazione
della libreria nativa (scritta in C) nell'ecosistema Python
di \texttt{HybridRepair} richiede una specifica gestione delle risorse.
Per abbattere il costo computazionale, il modulo implementa
un pattern architetturale noto come \textit{singleton lazy}. Di
seguito il frammento di codice Python incaricato di
questa orchestrazione:

\begin{lstlisting}[language=Python, label={lst:ensure_tree_sitter}]
def _ensure_tree_sitter():
    """Inizializzazione lazy del parser Tree-sitter come singleton."""
    global _TS_AVAILABLE, _JAVA_LANGUAGE, _PARSER
    
    if not _TS_AVAILABLE:
        from tree_sitter import Language, Parser
        import tree_sitter_java
        
        # Caricamento della grammatica compilata nativa
        _JAVA_LANGUAGE = Language(tree_sitter_java.language(), 'java')
        _PARSER = Parser()
        _PARSER.set_language(_JAVA_LANGUAGE)
        _TS_AVAILABLE = True
        
    return _PARSER
\end{lstlisting}

Il design pattern \textit{lazy singleton} interviene su un
collo di bottiglia fondamentale: l'inizializzazione da zero di
Tree-sitter (che prevede il caricamento a runtime della
grammatica pre-compilata tramite binding C) consuma svariate decine
di millisecondi. Al contrario, l'atto di analizzare e
derivare l'albero di un file (il parsing incrementale)
è estremamente reattivo e richiede solo una manciata
di millisecondi. Ritardando l'istanza al momento del primo
utilizzo e condividendo l'oggetto globale \texttt{\_PARSER} in memoria
per tutte le operazioni successive, il sistema previene
ritardi latenti deleteri durante l'esecuzione batch della suite
di riparazione. Il binding Python espone un puntatore
nativo, avvolto nell'oggetto \texttt{Language}, istruendo il parser a
operare secondo il dialetto specifico di Java.

\subsubsection{La Distinzione Cruciale: AST del Compilatore vs CST di Tree-sitter}
\label{subsubsec:3.4.5.3}
Un errore accademico comune è quello di interscambiare
i concetti di AST e CST. Comprendere il
salto qualitativo offerto da Tree-sitter richiede l'analisi della
dicotomia profonda tra l'astrazione semantica (AST) e la
fedeltà lessicale (CST).

\paragraph{L'Astrazione Semantica dell'AST del Compilatore}
Quando un compilatore convenzionale come \texttt{javac} ispeziona il
codice, il suo obiettivo ultimo è tradurre le
istruzioni in bytecode efficiente. Ne consegue che la
sua struttura ad albero, l'AST (\textit{Abstract Syntax Tree}),
si libera aggressivamente di tutti gli elementi superflui
che non impattano l'esecuzione del programma. Un AST
omette del tutto i commenti (siano essi \textit{block}
o \textit{line-comment}), ignora gli spazi bianchi e l'indentazione
visiva, rimuove il bilanciamento ridondante delle parentesi algebriche,
astrae i punti e virgola, e normalizza i
tipi letterali. Mentre questa spietata astrazione è inestimabile
per la compilazione, si rivela distruttiva per l'ingegneria
del software orientata alle patch. Se un agente
modificasse un nodo su un AST compilativo puro
e in seguito ritraducesse l'intero albero in codice
sorgente (processo di \textit{pretty-printing} o \textit{unparsing}), l'intero file
verrebbe ricostruito. Questa ricostruzione disgregherebbe irrimediabilmente l'indentazione e
lo stile originari adottati dagli ingegneri, riversando
modifiche estetiche su centinaia di righe non pertinenti
e rendendo impossibile la validazione del diff.

\paragraph{La Fedeltà Lessicale del CST di Tree-sitter}
In totale contrapposizione, Tree-sitter orchestra un Concrete Syntax
Tree (CST). Questa struttura non scarta nulla, vantando
una fedeltà lessicale \textit{pixel-perfect}. Ogni singolo commento, carattere
di spaziatura invisibile, ritorno a capo, annotazione e
parentesi trova il suo posto come nodo fisico
nell'albero. Il pregio definitivo del CST si
materializza negli attributi di posizionamento: ogni nodo trasporta
con sé uno \texttt{start\_byte} e un \texttt{end\_byte}, che
indicano gli offset precisi al byte all'interno del
flusso del file originale.

La rappresentazione ad albero sottostante illustra un frammento
minimale del CST restituito da Tree-sitter per una
dichiarazione di metodo Java:

\begin{verbatim}
method_declaration [byte: 0..46]
  modifiers: modifiers [byte: 0..6]
    public [byte: 0..6]
  type: void_type [byte: 7..11]
  name: identifier [byte: 12..16] "test"
  parameters: formal_parameters [byte: 16..18]
    ( [byte: 16..17]
    ) [byte: 17..18]
  body: block [byte: 19..46]
    { [byte: 19..20]
      ...
    } [byte: 45..46]
\end{verbatim}

Grazie agli indici di ancoraggio (es. \texttt{[byte: 0..46]}),
il sistema \texttt{HybridRepair} può estrarre, rimpiazzare o iniettare
blocchi logici operando unicamente in aritmetica sui byte,
con la certezza assoluta che lo \textit{slice} prelevato
in \texttt{source[node.start\_byte:node.end\_byte]} rappresenterà infallibilmente il testo fedele per
come è stato redatto nel sorgente originale, inclusivo
di tabulazioni e punteggiatura.

\subsubsection{Il Vocabolario dei Nodi Tree-sitter per Java}
\label{subsubsec:3.4.5.4}
La grammatica formale di \texttt{tree-sitter-java} è notevolmente vasta
e conta oltre 100 tipologie differenti di nodi
per descrivere ogni sfumatura del linguaggio. Nelle procedure
di sostituzione e astrazione intraprese da \texttt{HybridRepair}, vengono
selezionati in modo prevalente nodi in grado di
fornire i confini per le macro-strutture del codice.
Di seguito è riportata una tabella riassuntiva dei
principali nodi rilevanti:

\begin{table}[H]
\centering
\label{tab:treesitter_nodes}
\resizebox{\textwidth}{!}{
\begin{tabular}{l l l}
\toprule
\textbf{Tipo nodo} & \textbf{Descrizione} & \textbf{Campo figlio rilevante} \\
\midrule
\texttt{method\_declaration} & Dichiarazione di metodo & \texttt{name} (\texttt{identifier}), \texttt{body} (\texttt{block}) \\
\texttt{constructor\_declaration} & Dichiarazione di costruttore & \texttt{name} (\texttt{identifier}), \texttt{body} (\texttt{block}) \\
\texttt{field\_declaration} & Campo di classe & \texttt{type}, \texttt{declarator} \\
\texttt{local\_variable\_declaration} & Variabile locale & \texttt{type}, \texttt{declarator} \\
\texttt{formal\_parameter} & Parametro formale di metodo & \texttt{type}, \texttt{name} \\
\texttt{class\_declaration} & Dichiarazione di classe & \texttt{name}, \texttt{body} \\
\texttt{block} & Blocco di istruzioni & child nodes \\
\texttt{identifier} & Identificatore testuale Java & (testo del nodo) \\
\bottomrule
\end{tabular}
}
\end{table}

Questi nodi costituiscono il vocabolario attraverso cui il
sistema si orienta all'interno del file. A titolo
puramente illustrativo del modo in cui questi costrutti
informativi vengono impiegati dalla logica di riparazione di
\texttt{ast\_applier.py}, il seguente frammento espone l'approccio standard di
ispezione iterativa (tramite esplorazione a profondità DFS) per
localizzare il nodo esatto di un metodo target:

\begin{lstlisting}[language=Python, label={lst:find_method_by_name_intro}]
def _find_method_by_name(root_node, target_name):
    """Ricerca iterativa DFS del metodo target all'interno del CST."""
    stack = [root_node]
    
    while stack:
        node = stack.pop()
        
        # Intercettazione dichiarazioni di metodo o costruttore
        if node.type in ['method_declaration', 'constructor_declaration']:
            # Estrazione sicura del figlio denominato 'name'
            name_node = node.child_by_field_name('name')
            if name_node and name_node.text.decode('utf8') == target_name:
                return node
                
        # Aggiunta iterativa dei figli (in ordine decrescente 
        # per emulare una corretta visita left-to-right)
        for child in reversed(node.children):
            stack.append(child)
            
    return None
\end{lstlisting}

In questo breve passaggio algoritmico, il sistema naviga
fluidamente la gerarchia del CST, ispezionando unicamente i
nodi del tipo \texttt{method\_declaration} e \texttt{constructor\_dec\\laration} estratti dal
vocabolario, estraendo dal campo nominativo \texttt{name} (che corrisponde
all'\texttt{identifier} del metodo) il testo per il confronto
e ottenendo, infine, il nodo risolto portatore dei
corretti \textit{byte offset} su cui l'agente applicherà la
propria manipolazione strutturale.
\newpage
\subsection{Precisione Chirurgica e Modifiche Posizionali}
\label{subsec:3.4.6}
Il modulo \texttt{patch\_agent/ast\_applier.py} implementa il meccanismo centrale della
riparazione strutturale all'interno dell'ecosistema \texttt{HybridRepair}. La sua funzione
principale, \texttt{apply\_method\_replacement()}, ha il compito di concretizzare la
sostituzione posizionalmente precisa di un metodo nel sorgente Java.
Questa delicata operazione si articola in tre passaggi sequenziali:
il parsing del file attraverso Tree-sitter, la localizzazione inequivocabile
del nodo corrispondente al metodo target, e infine la
sostituzione del relativo span di testo. 

\subsubsection{L'Inizializzazione Lazy del Parser: \texttt{\_ensure\_tree\_sitter()}}
\label{subsubsec:3.4.6.1}
Il primo passo per poter operare sul codice è
l'istanza del parser. Come precedentemente introdotto, l'inizializzazione della libreria
avviene tramite la funzione \texttt{\_ensure\_tree\_s\\itter()}. Questa funzione non è
isolata, ma è condivisa in modo architetturalmente trasversale tra
i moduli \texttt{ast\_applier.py} (responsabile della riparazione) e \texttt{prolog\_groundi\\ng.py} (che
se ne avvale per la risoluzione dei tipi durante
la fase di \textit{grounding}). La condivisione sicura del singleton
del parser si basa sulle proprietà del sistema di
import di Python: le variabili globali del modulo, venendo
caricate una sola volta per processo, garantiscono che l'istanza
di Tree-sitter sia unica e condivisa in memoria tra
tutti gli importatori, abbattendo drasticamente i tempi di esecuzione
globale.

\subsubsection{La Funzione di Parsing: \texttt{\_parse\_java()}}
\label{subsubsec:3.4.6.2}
Una volta ottenuto il parser, il testo sorgente deve
essere analizzato per generare il \textit{Concrete Syntax Tree} (CST).
Questa operazione è delegata alla funzione \texttt{\_parse\_java()}, la quale
introduce un passaggio di codifica di fondamentale importanza:

\begin{lstlisting}[language=Python, label={lst:parse_java}]
def _parse_java(source_code: str):
    """Esegue il parsing del codice sorgente gestendo rigorosamente i byte."""
    parser = _ensure_tree_sitter()
    
    # Passaggio critico: codifica in bytes UTF-8
    source_bytes = source_code.encode('utf-8')
    
    # Generazione dell'albero sintattico
    tree = parser.parse(source_bytes)
    
    return tree, source_bytes
\end{lstlisting}

Il dettaglio implementativo critico visibile in questa funzione è
l'esplicita conversione della stringa Python in una sequenza di
byte tramite la codifica UTF-8 (\texttt{source\_code.en\\code('utf-8')}) prima di
invocare \texttt{parser.parse()}. Questa necessità deriva dall'architettura interna di
Tree-sitter, il quale opera esclusivamente nello spazio dei byte.
Di conseguenza, gli attributi posizionali restituiti dai nodi (come
\texttt{start\_byte} ed \texttt{end\_byte}) rappresentano offset in byte all'interno del
buffer codificato, non offset in caratteri Unicode. Mentre per
i file Java che contengono unicamente caratteri ASCII il
\textit{byte offset} e il \textit{character offset} coincidono perfettamente, la
situazione cambia radicalmente in presenza di stringhe o commenti
in lingue non-ASCII. Nei progetti internazionali, inclusi diversi applicativi
del dataset Defects4J sviluppati in Europa o in Asia,
l'utilizzo accidentale di \textit{character offset} provocherebbe disallineamenti silenti ma
letali durante la manipolazione, corrompendo la patch. Pertanto, operare
al livello dei byte UTF-8 è essenziale per la
sicurezza e la precisione dell'intero processo.

\subsubsection{La Strategia di Sostituzione via AST: \texttt{\_replace\_via\_ast()}}
\label{subsubsec:3.4.6.3}
Questa sezione descrive il cuore nevralgico del meccanismo di
riparazione strutturale. Una volta che il parser ha identificato
le coordinate esatte del metodo all'interno dell'albero, la sostituzione
fisica del codice viene eseguita dalla funzione \texttt{\_replace\_via\_ast()},
il cui codice è riportato di seguito:

\begin{lstlisting}[language=Python, label={lst:replace_via_ast}]
def _replace_via_ast(source_bytes: bytes, method_node, new_body_indented: str) -> bytes:
    """Esegue la sostituzione chirurgica sfruttando gli offset di Tree-sitter."""
    
    # Estrazione degli offset posizionali dal nodo
    start_byte = method_node.start_byte
    end_byte = method_node.end_byte
    
    # Codifica della patch generata dall'LLM
    new_body_bytes = new_body_indented.encode('utf-8')
    
    # Slicing posizionale e concatenazione
    patched_bytes = source_bytes[:start_byte] + new_body_bytes + source_bytes[end_byte:]
    
    return patched_bytes
\end{lstlisting}

La robustezza di questa funzione risiedono
nell'operazione di slicing e concatenazione \texttt{source\_bytes[:start\_byte] + new\_body\_bytes + source\_bytes[end\_byte:]}.
Analizziamo riga per riga la semantica di questa espressione:

\begin{description}
    \item[\texttt{source\_bytes[:start\_byte]}:] Questa porzione estrae in modo sicuro tutto il
    codice sorgente originale che precede esattamente il primo byte
    del metodo da sostituire, conservando intatti package, import, dichiarazioni
    di classe e metodi precedenti.
    
    \item[\texttt{new\_body\_bytes}:] Questa è la nuova implementazione corretta del metodo,
    generata dall'agente LLM e precedentemente re-indentata e codificata in
    UTF-8.
    
    \item[\texttt{source\_bytes[end\_byte:]}:] Questa istruzione recupera dal buffer originario tutto ciò
    che segue l'ultimo byte del metodo bersaglio (inclusa la
    parentesi graffa di chiusura della classe o i metodi
    successivi), completando il file.
\end{description}

L'affidabilità inossidabile di questa operazione è garantita da una
proprietà fondamentale di Tree-sitter: gli attributi \texttt{start\_byte} ed \texttt{end\_byte}
delineano in maniera assolutamente esatta e priva di ambiguità
posizionale l'intervallo del file sorgente originale di competenza del
nodo. Da questa meccanica deriva un teorema pratico di
correttezza sintattica per costruzione. La proprietà formale asserisce che:
se il file \texttt{source\_bytes} originale è sintatticamente valido, se
\texttt{new\_body\_indented} è a sua volta un metodo sintatticamente valido
con il medesimo identificatore, e se \texttt{method\_node} inquadra correttamente
le coordinate nel CST, allora l'output \texttt{patched\_bytes} risultante è
matematicamente garantito essere un file Java sintatticamente integro. A
differenza degli approcci testuali basati su diff o regex,
dove la validità del risultato è precaria e dipendente
dalla forma esteriore del testo, qui la correttezza non
è una scommessa statistica, ma una conseguenza diretta della
struttura formale della manipolazione posizionale. Eventuali errori logici o
semantici residui nella patch verranno poi demandati e intercettati
dal compilatore nella successiva fase \texttt{COMPILE}.

\subsubsection{Il Rilevamento del Metodo Target: \texttt{\_find\_method\_by\_name()}}
\label{subsubsec:3.4.6.4}
Per poter alimentare la funzione di sostituzione posizionale con
il nodo corretto, il sistema deve prima localizzarlo all'interno
del CST. Questa ricerca è implementata dalla funzione \texttt{\_find\_method\_by\_name()},
che incorpora una logica complessa per gestire gerarchie profonde
e risolvere ambiguità:

\begin{lstlisting}[language=Python, label={lst:find_method_advanced}]
def _find_method_by_name(root_node, target_name: str, target_line: int = None):
    """Ricerca iterativa DFS del metodo target con disambiguazione degli overload."""
    stack = [root_node]
    candidate_nodes = []
    
    while stack:
        node = stack.pop()
        
        # Identificazione di metodi o costruttori
        if node.type in ['method_declaration', 'constructor_declaration']:
            name_node = node.child_by_field_name('name')
            if name_node and name_node.text.decode('utf-8') == target_name:
                candidate_nodes.append(node)
                
        # Inserimento inverso dei figli per preservare la visita left-to-right
        for child in reversed(node.children):
            stack.append(child)
            
    # Disambiguazione in presenza di overload multipli
    if target_line is not None and len(candidate_nodes) > 1:
        for node in candidate_nodes:
            # Tree-sitter restituisce righe zero-indexed, aggiungiamo 1
            start_row = node.start_point + 1
            end_row = node.end_point + 1
            
            # Verifica se la linea difettosa ricade nei confini del metodo
            if start_row <= target_line <= end_row:
                return node
                
    return candidate_nodes if candidate_nodes else None
\end{lstlisting}

Ingegneristicamente, due aspetti di questa implementazione spiccano per importanza.
In primo luogo, l'esplorazione del CST avviene tramite un
algoritmo \textit{Depth-First Search} (DFS) rigorosamente iterativo, supportato da uno
stack esplicito, ripudiando il più canonico approccio ricorsivo. Questa
è una scelta mirata alla robustezza strutturale, resa necessaria
dalle caratteristiche dei file Java di livello \textit{enterprise}. Nei
progetti industriali, l'annidamento sintattico può raggiungere profondità estreme (ad
esempio, a causa di classi anonime intrecciate con blocchi
\texttt{try/catch/finally} e iterazioni annidate). Un DFS ricorsivo rischierebbe di
superare rapidamente il limite prefissato dello stack di chiamate
del runtime Python (tipicamente limitato a 1000 frame), scatenando
un fatale \texttt{RecursionError}. Lo stack iterativo, gestito dinamicamente nell'heap
di memoria, immunizza l'algoritmo da questa specifica vulnerabilità.

In secondo luogo, la funzione implementa un meccanismo
di disambiguazione guidato dal parametro \texttt{target\_line}. Nel vasto ecosistema
di Defects4J, classi cardine come \texttt{DatasetUtilities} (nel progetto JFreeChart)
o \texttt{StringUtils} (in Apache Commons Lang) presentano frequentemente dozzine
di metodi sovraccaricati (\textit{overload}), recanti il medesimo nome identificativo
ma differenti firme e comportamenti. Se il sistema rintracciasse
semplicemente il nome (ottenendo molteplici nodi candidati in \texttt{candidate\_nodes}),
si affiderebbe al caso applicando la patch al primo
della lista. Ricevendo la \texttt{target\_line}, che il \texttt{PatchAgent}
propaga attingendo alle informazioni di \texttt{FaultLocation} presenti nel \texttt{FaultReport},
 l'algoritmo confronta questa coordinata con il \textit{bounding box}
di ogni metodo candidato (\texttt{start\_point} ed \texttt{end\_point}). Solo il
nodo che abbraccia fisicamente la riga difettosa viene selezionato,
risolvendo in modo univoco l'ambiguità degli overload
multipli.

\subsubsection{La Preservazione dell'Indentazione: \texttt{\_detect\_indentation()} e \texttt{\_reindent()}}
\label{subsubsec:3.4.6.5}
La preservazione dello stile di indentazione originale del codice
 costituisce un
requisito funzionale critico per la correttezza pratica del sistema.
Se il \texttt{PatchAgent} iniettasse un metodo generato dall'LLM con
un'indentazione discorde rispetto a quella del file target, il
sistema produrrebbe un diff ``inquinato'', in cui ogni singola
riga del metodo risulterebbe modificata a causa delle discrepanze
negli spazi bianchi, e non solo le righe che
hanno effettivamente subito un'alterazione logica. Questo disallineamento renderebbe la
patch estremamente difficile da revisionare per un operatore umano
e aumenterebbe drasticamente la probabilità di introdurre \textit{whitespace-related compilation
errors}, specialmente in sistemi di \textit{Continuous Integration} (CI) che
impongono regole di stile stringenti (come Checkstyle).

Per scongiurare questo problema, il sistema implementa due funzioni
dedicate: \texttt{\_detect\_\\indentation()} per campionare lo stile originario e \texttt{\_reindent()}
per applicarlo dinamicamente al codice generato.

\begin{lstlisting}[language=Python, label={lst:detect_indentation}]
def _detect_indentation(source_code: str) -> str:
    """Rileva l'indentazione del prefisso della prima riga utile."""
    first_line = source_code.split('\n')[0]
    # Estrazione del prefisso composto esclusivamente da spazi o tab
    indentation_len = len(first_line) - len(first_line.lstrip())
    return first_line[:indentation_len]
\end{lstlisting}

La funzione \texttt{\_detect\_indentation()} si occupa di estrarre il prefisso
esatto (spazi bianchi ASCII o tabulazioni) della prima riga
del metodo originale nel file target. Nel dataset Defects4J,
le convenzioni stilistiche presentano un'elevata variabilità: progetti come JFreeChart
adottano rigorosamente 4 spazi, alcuni sottoprogetti di Apache Commons
privilegiano le tabulazioni, mentre Closure Compiler impiega tipicamente 2
spazi. Il rilevamento dinamico svincola il sistema dalla necessità
di conoscere a priori le policy del progetto bersaglio.

Una volta estratto il prefisso, il codice sostitutivo fornito
dall'agente deve essere riallineato tramite \texttt{\_reindent()}:

\begin{lstlisting}[language=Python, label={lst:reindent}]
def _reindent(new_body: str, target_indent: str) -> str:
    """Applica una trasformazione di indentazione relativa al nuovo corpo del metodo."""
    lines = new_body.split('\n')
    if not lines:
        return new_body
        
    # Calcolo dell'indentazione base (livello esterno) del codice LLM
    base_indent_len = len(lines[0]) - len(lines[0].lstrip())
    base_indent = lines[0][:base_indent_len]
    
    reindented_lines = []
    for line in lines:
        # Gestione esplicita delle righe vuote
        if not line.strip():
            reindented_lines.append("")
            continue
            
        # Trasformazione relativa: sostituisce l'indentazione base con quella target
        if line.startswith(base_indent):
            reindented_lines.append(target_indent + line[base_indent_len:])
        else:
            reindented_lines.append(target_indent + line.lstrip())
            
    return '\n'.join(reindented_lines)
\end{lstlisting}

Un aspetto fondamentale è la differenza tra un'imposizione
di indentazione assoluta e la trasformazione di indentazione relativa
qui implementata. Un approccio assoluto appiattirebbe forzatamente ogni riga
a un margine fisso, distruggendo la struttura visiva del
codice. L'algoritmo relativo, invece, calcola il delta tra l'indentazione
esterna del codice generato dall'LLM e il \texttt{target\_indent} prelevato
dal file fisico; questa differenza viene poi applicata uniformemente
a ogni riga. In questo modo, l'indentazione interna del
metodo, ovvero i delicati livelli di annidamento relativi
di blocchi \texttt{if}, cicli \texttt{for} e costrutti \texttt{try/catch},
viene perfettamente preservata, traslando semplicemente l'intero blocco logico al
livello di rientro corretto del file ospite. Da notare
anche la gestione esplicita delle righe vuote (\texttt{if not line.strip()}).
Questa riga di codice è necessaria per preservare intatte
le spaziature verticali di separazione logica introdotte all'interno del
metodo, che nei linguaggi ad alta leggibilità come Java
costituiscono convenzioni visive fondamentali per la manutenibilità.

\subsubsection{Il Calcolo del Diff Unificato}
\label{subsubsec:3.4.6.6}
Sebbene il sistema abbia definitivamente abbandonato i diff come
veicolo per applicare la modifica, il formato \textit{unidiff} mantiene
un ruolo di primaria importanza a livello diagnostico.

\begin{lstlisting}[language=Python, label={lst:compute_diff}]
def _compute_diff(original: str, patched: str, filename: str) -> str:
    """Genera un diff unificato diagnostico post-applicazione."""
    import difflib
    
    original_lines = original.splitlines(keepends=True)
    patched_lines = patched.splitlines(keepends=True)
    
    # Parametro n=3 per rispettare lo standard delle righe di contesto
    diff = difflib.unified_diff(
        original_lines, patched_lines,
        fromfile=filename, tofile=filename, n=3
    )
    return "".join(diff)
\end{lstlisting}

La funzione \texttt{\_compute\_diff()} viene rigorosamente invocata solo dopo l'applicazione
della patch tramite Tree-sitter, assumendo un valore puramente diagnostico
e visivo per il modello linguistico. Il risultato non
viene mai utilizzato per alterare il file sorgente; al
contrario, viene formattato e restituito all'agente LLM come feedback
di conferma (\texttt{"[SUCCESS] Patch applied successfully ... Diff: ..."})
per permettergli di validare visivamente l'impatto della sua azione.
Il parametro \texttt{n=3} ricalca deliberatamente la convenzione standard di
\texttt{git diff} per le righe di contesto: valori superiori
aumenterebbero la leggibilità a uso umano, ma si tradurrebbero
in stringhe troppo verbose che consumerebbero inutilmente il budget
di token contestuali (\textit{token budget}) della finestra di memoria
dell'LLM.

\subsubsection{La Strategia di Fallback Graduale}
\label{subsubsec:3.4.6.7}
L'orchestrazione finale del processo di sostituzione strutturale è racchiusa
nel metodo \texttt{apply\_method\_replacement()}. Questa funzione si fa carico di
integrare il parsing, l'indentazione, la manipolazione dell'AST e la
generazione del diff in un'unica interfaccia robusta per il
\texttt{PatchAgent}.

Per garantire la massima resilienza, specialmente considerando l'eterogeneità di
ambienti in cui il software potrebbe essere eseguito, il
metodo implementa una strategia di fallback graduale a
due livelli:

\newpage
\begin{lstlisting}[language=Python, label={lst:apply_method_replacement}]
def apply_method_replacement(file_path: Path, method_name: str, new_body: str, target_line: int = None) -> tuple[bool, str]:
    """Applica la sostituzione del metodo tramite una strategia di fallback graduale."""
    original_source = file_path.read_text(encoding='utf-8')
    
    # 1. Preservazione stile
    target_indent = _detect_indentation(original_source)
    new_body_indented = _reindent(new_body, target_indent)
    
    # 2. LIVELLO PRIMARIO: Manipolazione AST via Tree-sitter
    try:
        tree, source_bytes = _parse_java(original_source)
        method_node = _find_method_by_name(tree.root_node, method_name, target_line)
        
        if method_node:
            patched_bytes = _replace_via_ast(source_bytes, method_node, new_body_indented)
            patched_source = patched_bytes.decode('utf-8')
            
            file_path.write_text(patched_source, encoding='utf-8')
            diff = _compute_diff(original_source, patched_source, file_path.name)
            return True, f"[SUCCESS] Patch applied successfully.\nDiff:\n{diff}"
    except Exception as e:
        # Errore non fatale, si procede al livello di fallback
        pass

    # 3. LIVELLO SECONDARIO: Fallback Regex
    try:
        # L'approccio regex richiede obbligatoriamente il target_line
        patched_source = _apply_regex_fallback(original_source, method_name, new_body_indented, target_line)
        
        if patched_source:
            file_path.write_text(patched_source, encoding='utf-8')
            diff = _compute_diff(original_source, patched_source, file_path.name)
            return True, f"[WARNING] Patch applied via Regex fallback.\nDiff:\n{diff}"
    except Exception as e:
        pass
        
    # 4. LIVELLO FINALE: Fallimento Graceful
    return False, "[FAILED] Failed to apply patch: method boundaries not found or parsing failed."
\end{lstlisting}

Questa ridondanza architetturale definisce un ordine di priorità strettamente
decrescente in termini di affidabilità. Il Livello Primario tenta
sistematicamente la sostituzione via Tree-sitter AST, che garantisce la
correttezza sintattica per costruzione, la gestione nativa degli overload
e la compatibilità con annotazioni e gerarchie di annidamento
arbitrarie. Qualora il parsing dovesse fallire (circostanza rara, ma
possibile in caso di file sorgente severamente corrotto in
partenza) o l'infrastruttura binaria nativa in C della libreria
\texttt{tree-sitter-java} non risultasse accessibile nel sistema operativo in uso,
subentra il Livello Secondario. Questo si affida a un
approccio basato su espressioni regolari Il fallback regex opera
solo su file Java standard, richiede obbligatoriamente il parametro
\texttt{target\_line} e, seppur soggetto al rischio potenziale di sbilanciare
parentesi in metodi altamente complessi, offre una via di
fuga operazionale vitale in ambienti non standardizzati privi di
dipendenze binarie. 

Se anche l'estrazione testuale fallisce clamorosamente, il sistema invoca
il Fallimento Graceful. Invece di interrompere brutalmente l'esecuzione sollevando
un'eccezione non gestita a livello di thread che causerebbe
il crash dell'agente, il sistema restituisce in modo sicuro
una tupla di errore testuale. Questo messaggio negativo viene
iniettato nel contesto dell'LLM; l'agente riceve notifica del fallimento
materiale e può optare, nel round successivo, per un
approccio logico differente o fornire indizi addizionali per aiutare
il parser. Questa \textit{graceful degradation} è un principio cardine
dell'intera architettura, risultando imprescindibile per un sistema APR automatizzato
chiamato a processare batch massivi e ininterrotti di bug
profondamente eterogenei.

\subsection{Il Sandbox di Valutazione: \texttt{sandbox\_evaluator.py}}
\label{subsec:3.4.7}
Il modulo \texttt{services/sandbox\_evaluator.py} implementa il sandbox di valutazione,
l'ambiente isolato che fornisce all'agente i due oracoli deterministici
necessari per il ciclo della Macchina a Stati Finiti:
il compilatore Java (\texttt{javac}) e il framework di esecuzione
dei test (\texttt{org.junit.runner.JUnitCore}). In ambito di ingegneria del software
orientata alla sicurezza, il termine ``sandbox'' denota un rigoroso
isolamento operativo. Ogni singolo tentativo di riparazione agisce su
una copia fisica separata del codice sorgente, confinata in
una directory dedicata, per garantire che i fallimenti e
le mutazioni di un tentativo non inquinino lo stato
dei successivi, mantenendoli reciprocamente indipendenti.

\subsubsection{La Gestione del Classpath: \texttt{config.properties} e Glob Expansion}
\label{subsubsec:3.4.7.1}
Prima di poter compilare o eseguire il codice Java,
l'ambiente necessita di una corretta risoluzione delle dipendenze. Nel
dataset Defects4J, questa configurazione è memorizzata all'interno di un
file specifico, \texttt{config.properties}, presente nella directory radice di ogni
bug. L'estrazione e la risoluzione di questo classpath avvengono
tramite la funzione \texttt{\_read\_classpath()}:

\begin{lstlisting}[language=Python, label={lst:read_classpath}]
def _read_classpath(bug_dir: Path) -> str:
    """Estrae e risolve il classpath dal file config.properties."""
    config_path = bug_dir / "config.properties"
    
    # Estrazione della riga contenente class.path
    with open(config_path, 'r') as f:
        for line in f:
            if line.startswith("class.path="):
                raw_paths = line.strip().split("=")[3].split(":")
                break
                
    resolved_entries = []
    for entry in raw_paths:
        if '*' in entry:
            # Meccanismo di glob expansion per le wildcard
            base_dir = (bug_dir / entry).parent
            for jar_file in base_dir.glob("*.jar"):
                resolved_entries.append(str(jar_file.resolve()))
        else:
            # Risoluzione in percorso assoluto
            resolved_entries.append(str((bug_dir / entry).resolve()))
            
    return ":".join(resolved_entries)
\end{lstlisting}

Il file \texttt{config.properties} definisce il classpath utilizzando percorsi relativi alla
directory del bug, concatenandoli tipicamente in una forma del
tipo \texttt{class.path=../../lib/j\\unit-4.11.jar:../../lib/hamcrest-core-1.3.jar:...}. La funzione si occupa di ancorare
queste definizioni relative trasformandole in percorsi assoluti tramite l'invocazione
di \texttt{(bug\_dir / entry).resolve()}. Un aspetto architetturale di vitale importanza
è il supporto per la \textit{glob expansion}. Alcuni progetti
non elencano minuziosamente ogni singolo archivio JAR, ma utilizzano
sintassi con wildcard (come \texttt{../../lib/*}). Per intercettare questa notazione,
il codice verifica la presenza del carattere \texttt{*} e
innesca la libreria \texttt{glob}, iterando il filesystem per raccogliere
dinamicamente tutti gli archivi \texttt{.jar} presenti nella directory. Questo
meccanismo astratto assorbe con eleganza le marcate differenze strutturali
nella gestione delle dipendenze dei vari repository che compongono
Defects4J.

\subsubsection{La Compilazione Selettiva dei File Modificati}
\label{subsubsec:3.4.7.2}
La verifica sintattica del codice modificato richiede l'uso del
compilatore \texttt{javac}. Tuttavia, ricompilare un intero progetto enterprise (come
JFreeChart o Closure Compiler) per ogni singolo tentativo della
FSM richiederebbe interi minuti, rendendo proibitivi i tempi di
esecuzione del sistema. Il sandbox supera questa limitazione implementando
un approccio di compilazione selettiva:

\begin{lstlisting}[language=Python, label={lst:compile_patched_source}]
def _find_modified_files(original_dir: Path, patched_dir: Path) -> list[str]:
    """Identifica i file Java modificati tramite confronto byte-a-byte."""
    import filecmp
    comparison = filecmp.dircmp(original_dir, patched_dir)
    return [f for f in comparison.diff_files if f.endswith('.java')]

def compile_patched_source(patched_dir: Path, bug_dir: Path, classpath: str):
    """Compila selettivamente il file modificato tramite javac."""
    import subprocess
    
    modified_files = _find_modified_files(bug_dir, patched_dir)
    if not modified_files:
        return True, ""
        
    target_file = patched_dir / modified_files[0]
    
    cmd = [
        "javac",
        "-cp", classpath,
        "-sourcepath", str(patched_dir),
        "-nowarn",
        "-encoding", "UTF-8",
        str(target_file)
    ]
    
    result = subprocess.run(cmd, capture_output=True, text=True)
    return result.returncode == 0, result.stderr
\end{lstlisting}

La funzione sussidiaria \texttt{\_find\_modified\_files()} impiega il modulo \texttt{filecmp.dircmp()} per
operare un confronto byte-a-byte tra la directory originale e
il sandbox. Questa scelta è strategicamente superiore a un
banale controllo dei timestamp: la copia dei file tra
file system potrebbe inavvertitamente aggiornare i timestamp, innescando falsi
positivi, mentre il confronto binario estrae infallibilmente il solo
file che ha subito mutazioni da parte dell'LLM.

Una volta isolato il file (es. \texttt{PatchedClass.java}), viene invocato
il compilatore. Le opzioni passate a \texttt{javac} costituiscono il
fulcro concettuale di questa procedura. Il flag \texttt{-sourcepath} è
di importanza primordiale: esso istruisce il compilatore a risolvere
i riferimenti a classi non specificate sulla riga di
comando andandone a cercare i sorgenti all'interno della directory
radice indicata. Grazie a questo parametro, \texttt{javac} può compilare
unicamente la \texttt{PatchedClass.java} modificata, riuscendo contemporaneamente a risolvere i
riferimenti incrociati alle altre classi dello stesso package, scongiurando
i costi di una ricompilazione totale del progetto. Parimenti
cruciale è l'inclusione del flag \texttt{-nowarn}. Il codice legacy
abbonda di pratiche superate che generano avvisi (come \textit{deprecated
method calls} o \textit{unchecked assignments} sui generici). Questi messaggi non
indicano violazioni bloccanti e non interrompono la compilazione, ma,
se passati all'agente LLM, inquinerebbero il prompt con rumore
superfluo. Omettendoli, ci si assicura che il modello si
concentri esclusivamente sui veri errori sintattici.
\newpage
\subsubsection{L'Esecuzione JUnit e il Parsing dell'Output}
\label{subsubsec:3.4.7.3}
L'oracolo semantico finale del ciclo di riparazione è demandato
al collaudo della suite di test. Poiché il sistema
opera in totale automazione e disaccoppiamento dagli IDE grafici,
i test vengono innescati invocando \texttt{org.junit.runner.JUnitCore} in pura modalità
\textit{command-line}:

\begin{lstlisting}[language=Python,  label={lst:run_junit}]
def _run_junit(patched_dir: Path, classpath: str, test_classes: list[str]):
    """Esegue la suite JUnit in modalità command-line con limite temporale."""
    import subprocess
    
    cmd = [
        "java", 
        "-cp", classpath, 
        "org.junit.runner.JUnitCore"
    ] + test_classes
    
    # Timeout di 300 secondi per prevenire stalli
    return subprocess.run(cmd, capture_output=True, text=True, timeout=300)
\end{lstlisting}

In questa esecuzione, il parametro \texttt{timeout=300} svolge il ruolo
di scudo  contro un fenomeno  dell'Automated Program
Repair: le modifiche introdotte dall'agente, specialmente su bug legati
alla sincronizzazione multithreading o ai flussi di I/O, possono
accidentalmente generare deadlock o loop infiniti (\textit{infinite loops}) all'interno
dei test. Fissando un tempo limite di 5 minuti,
si impone un vincolo di sicurezza che previene il
blocco indefinito dell'intero processo di elaborazione batch.

Il runner JUnit standard in modalità CLI non restituisce
un log strutturato (come un XML), bensì riversa sullo
standard output un report testuale grezzo. Questo output deve
essere imperativamente analizzato per estrarre logiche diagnostiche avanzate. L'operazione
è affidata alla funzione \texttt{\_parse\_junit\_output()}:

\begin{lstlisting}[language=Python , label={lst:parse_junit_output}]
def _parse_junit_output(output: str, original_failing_tests: set):
    """Analizza l'output testuale JUnit ed estrae le metriche di progresso."""
    import re
    
    # Estrazione delle firme tramite espressione regolare
    failure_pattern = re.compile(r"\.E\s+([a-zA-Z0-9_]+)\(([a-zA-Z0-9_.]+)\)")
    
    current_failures = set()
    for match in failure_pattern.finditer(output):
        method_name, class_name = match.groups()
        current_failures.add(f"{class_name}::{method_name}")
        
    previously_failing = original_failing_tests
    still_failing = previously_failing.intersection(current_failures)
    
    # Metrica cruciale di parziale successo
    previously_failing_now_passing = previously_failing - still_failing
    
    return {
        "status": "PASS" if not current_failures and "OK (" in output else "TEST_FAIL",
        "failures": list(current_failures),
        "previously_failing_now_passing": list(previously_failing_now_passing)
    }
\end{lstlisting}

L'algoritmo impiega un'espressione regolare (\verb|\.E\s+...|) per ispezionare
il log testuale ed estrarre le firme esatte dei
test falliti (scorporando il nome del metodo da quello
della classe originaria). L'elenco ottenuto (\texttt{current\_failures}) viene incrociato dinamicamente
con l'\texttt{original\_failing\_tests} (la mappa dei fallimenti pre-patch contenuta nel
\texttt{FaultReport}) per ricavare l'attributo \texttt{previously\_failing\_now\_passing}. Il calcolo di questa
metrica costituisce un apporto metodologico di straordinaria utilità. In
molti frangenti di riparazione automatica complessa, una singola patch
non riesce a far transitare tutti i test in
stato verde simultaneamente. Misurare con esattezza quanti e quali
test specifici abbiano invertito il loro stato da fallito
a corretto (ad esempio, registrando 3 test passati su
5) offre un indicatore di progresso.
Esso segnala all'agente che la rotta esplorativa è parzialmente
valida e promettente, fornendo un fondamentale indicatore direzionale pur
in assenza del raggiungimento totale del criterio di successo
formale.

\subsection{Strutture Dati della Fase 3: \texttt{PatchAttempt} e \texttt{RepairResult}}
\label{subsec:3.4.8}
Il contratto di interfaccia e di output della Fase
3 è rigorosamente formalizzato attraverso due strutture dati immutabili,
definite come \texttt{@dataclass} all'interno del modulo \texttt{shared/models.py}. L'ingegnerizzazione
di queste classi riflette una precisa filosofia di design:
esse non si limitano a registrare il mero esito
binario (successo o fallimento) dell'operazione di riparazione, ma fungono
da veri e propri archivi permanenti. Di
seguito viene presentata la definizione architetturale di queste due
entità:

\begin{lstlisting}[language=Python, label={lst:models_patch_repair}]
from dataclasses import dataclass, field
from typing import Optional, List, Dict, Any

@dataclass
class PatchAttempt:
    """Archivio permanente di un singolo tentativo di riparazione."""
    id: int
    success: bool
    agent_messages: List[Dict[str, Any]] = field(default_factory=list)
    applied_diff: Optional[str] = None
    latest_error: Optional[str] = None
    test_results: Optional[Dict[str, Any]] = None

@dataclass
class RepairResult:
    """Artefatto finale aggregato dell'intera pipeline HybridRepair."""
    bug_id: str
    success: bool
    winning_attempt: Optional[int] = None
    attempts: List[PatchAttempt] = field(default_factory=list)
    fault_report: Any = None  # Riferimento alla struttura FaultReport
    repair_spec: Any = None   # Riferimento alla struttura RepairSpec
    
    def summary(self) -> str:
        """Produce una stringa riassuntiva per l'output a terminale."""
        status = "RESOLVED" if self.success else "FAILED"
        winning = f" (Attempt {self.winning_attempt})" if self.success else ""
        return f"Bug {self.bug_id}: {status}{winning} - Total attempts: {len(self.attempts)}"
\end{lstlisting}

La struttura \texttt{PatchAttempt} è stata  progettata per abilitare
una completa tracciabilità post-hoc. Il suo attributo più gravido
di informazioni è \texttt{agent\_messages}, il quale incapsula l'intera storia
conversazionale intercorsa tra l'agente LLM e il sistema durante
quello specifico tentativo. Questo include ogni singolo prompt di
sistema, le invocazioni strutturate agli strumenti (con i relativi
identificatori \texttt{tool\_call\_id}), i diff testuali generati a scopo diagnostico,
gli output troncati del compilatore e i fallimenti catturati
da JUnit. Questa granularità permette di ricostruire a posteriori,
passo dopo passo, l'intera catena decisionale e deduttiva del
modello generativo, rivelandosi uno strumento insostituibile per il debugging
algoritmico e l'analisi del comportamento emergente dell'agente.

Ad un livello di astrazione superiore si colloca \texttt{RepairResult},
che costituisce l'artefatto conclusivo e aggregato dell'intera pipeline a
tre fasi. Questa struttura raccoglie al suo interno non
solo la lista sequenziale di tutti i \texttt{PatchAttempt} eseguiti,
ma ingloba anche il \texttt{FaultReport} (ereditato dalla Fase 1)
e la \texttt{RepairSpec} (elaborata nella Fase 2), unificando così
diagnosi, specifica e risoluzione in un unico blocco informativo.
Il metodo \texttt{summary()} fornisce un'astrazione concisa, emettendo una singola
riga di testo che si rivela essenziale per il
monitoraggio a terminale durante l'esecuzione batch e concorrente su
centinaia di bug (tramite lo script \texttt{run\_all\_bugs.py}).

Al termine dell'esecuzione, l'onere della persistenza è delegato al
modulo \texttt{services/rep\\orter.py}, il quale serializza ricorsivamente l'oggetto \texttt{RepairResult} in
formato JSON. Il file risultante, salvato sotto il nome
di \texttt{final\_report.json} all'interno della directory isolata \texttt{pipeline\_results/<bug\_id>/}, si cristallizza
in un archivio storico permanente. Questa persistenza documentale è
inestimabile per la ricerca accademica, poiché fornisce il dataset
grezzo su cui condurre successive analisi statistiche di efficacia
e comparazioni rigorose con altri sistemi di Automated Program
Repair presenti in letteratura.

\subsection{Il Flusso Completo: dalla \texttt{RepairSpec} al \texttt{RepairResult}}
\label{subsec:3.4.9}
Per sintetizzare la vasta orchestrazione tecnica fin qui delineata,
è possibile tracciare l'intero viaggio dei dati all'interno del
\texttt{PatchAgent}. Questo flusso algoritmico, che unisce inferenza neurale e
oracoli deterministici, si snoda attraverso tre passaggi fondamentali:

\paragraph{Passo 1 --- Inizializzazione del Contesto Strategico}
Il ciclo di vita prende avvio nel costruttore \texttt{PatchAgent.\_\_init\_\_()}.
L'agente riceve in ingresso il trittico di artefatti che
fonda la sua conoscenza del problema: il \texttt{FaultReport} (la
mappa del difetto), la \texttt{RepairSpec} (la strategia risolutiva) e
gli \texttt{Ingredients} (il vocabolario delle API ammesse). In questa
fase preparatoria, viene istanziato il client di comunicazione \texttt{AzureClient}
e viene pre-allocato un oggetto \texttt{RepairResult} con lo stato
iniziale impostato rigorosamente a \texttt{success=False}.

\paragraph{Passo 2 --- Il Loop della Macchina a Stati Finiti (Max 3 Tentativi)}
Il cuore operativo dell'architettura innesca un ciclo iterativo che
concede all'agente un massimo di tre tentativi sequenziali (\texttt{max\_attempts = 3})
per giungere a una soluzione. Per ogni tentativo $k$:

\begin{description}
    \item[Isolamento Operativo:] La funzione \texttt{\_create\_patched\_copy()} clona il progetto target
    in una directory sandbox vergine (\texttt{pipeline\_results/<bug\_id>/attempt\_k/p\\atched\_source/}), azzerando il rischio
    di contaminazioni inter-tentativo.
    
    \item[Modulazione Stocastica:] Viene prelevata la temperatura di decodifica $T_k$
    dal vettore \texttt{LLM\_TEMPERATURE\_SCHEDULE} (valori 0.1, 0.2, 0.4), aumentando progressivamente
    il parametro di ``creatività'' del modello a fronte dei
    pregressi fallimenti.
    
    \item[Innesco del Tool-Use Loop:] Viene invocata la routine \texttt{\_run\_single\_attempt()}.
    Si istanzia il \texttt{ToolContext}, si predispone il dispatcher degli
    strumenti (\texttt{create\_tool\_\\executor()}) e si crea il prompt iniziale tramite
    \texttt{\_build\_messages()}. Se l'agente si trova in un tentativo
    $k > 1$, il prompt viene arricchito iniettando i
    riassunti dei fallimenti precedenti estratti dalla memoria conversazionale.
    
    \item[Dinamica Conversazionale:] Il controllo passa ad \texttt{AzureClient.chat\_with\_tools()}, che avvia
    un loop autonomo vincolato da un tetto massimo di
    9 iterazioni (\texttt{max\_too\\l\_rounds}). Durante ogni round, il modello si muove
    tra gli stati della FSM (\texttt{EXPLORE}, \texttt{PATCH}, \texttt{COMPILE},
    \texttt{TEST}), invocando strumenti, manipolando l'albero CST del codice e
    analizzando gli output generati. Questo anello si spezza unicamente
    quando il modello emette un responso testuale finale privo
    di \texttt{tool\_calls} o quando esaurisce il budget di round.
    
    \item[Persistenza ed Esame:] Conclusa l'interazione generativa, \texttt{conversation\_store.sav\\e()} interviene per
    serializzare l'intero scambio in un log JSONL. Immediatamente
    dopo, il giudice supremo \texttt{\_verify\_final()} entra in scena, invocando
    congiuntamente il compilatore \texttt{javac} e il runner JUnit sulla
    sandbox per produrre un \texttt{TestResult} inappellabile.
\enlargethispage{2\baselineskip}

    \item[Valutazione del Verdetto:] Se l'oracolo restituisce \texttt{PASS}, il traguardo è
    raggiunto: \texttt{RepairResult.success} diventa \texttt{True}, il numero del tentativo corrente
    viene marcato come \texttt{winning\_attempt} e il macro-ciclo viene interrotto
    prematuramente con successo. Se il verdetto differisce, l'iterazione fallisce
    e l'agente sfocia nel tentativo successivo, forte di una
    nuova lezione memorizzata.
\end{description}

\paragraph{Passo 3 --- Restituzione del \texttt{RepairResult}}
Esauriti i tentativi a disposizione o conseguito il successo,
il viaggio giunge al termine. L'oggetto \texttt{RepairResult}, ormai ricolmo
di dati diagnostici, diff, messaggi e metriche, viene restituito
per intero all'orchestratore di livello superiore \texttt{repair\_bug.py}. Quest'ultimo, richiamando
\texttt{reporter.py}, imprime definitivamente l'oggetto su disco, sigillando la riparazione
del bug e liberando le risorse computazionali per l'analisi
del difetto successivo.


